\section{Method}
\label{sec:method}

\subsection{Setup and notation}
\label{sec:notation}
Let $Q$ be a set of questions and $C$ a fixed passage corpus.
A retrieval method $R$ maps a question $q \in Q$ to a ranked list of passages from $C$.
A reader $M$ receives the question and the top $k$ passages and returns a short answer $a$.
A judge $J$ receives the same passages, the question and $a$, and returns SUPPORTED or NOT\_SUPPORTED.
Each question $q$ has a set of relevant passages $G(q) \subset C$ from the BEIR relevance judgements \citep{thakur2021beir} and a set of reference answers from NQ-open.

We use three per-query measures.
Hit@$k$ is 1 if at least one passage in $G(q)$ appears in the top $k$ of $R(q)$, and 0 otherwise.
nDCG@10 is the normalised discounted cumulative gain of the top 10 passages under the BEIR judgements.
Exact match (EM) is 1 if the answer equals a reference answer after SQuAD-style normalisation (lower-casing and removal of punctuation, articles and extra whitespace), and 0 otherwise.
Each reported score is the mean of a per-query measure over the query set.

\subsection{Transfer decomposition}
\label{sec:transfer}
Our central quantity compares two retrieval methods on the same queries.
For a method $R$ and the reference $R_1$ (all-MiniLM-L6-v2), we define
\begin{equation}
T(R) = \frac{\mathrm{EM}(R) - \mathrm{EM}(R_1)}{\mathrm{Hit@5}(R) - \mathrm{Hit@5}(R_1)} .
\label{eq:transfer}
\end{equation}
$T(R)$ is the change in EM per unit change in Hit@5.
A value of 1 means that every additional retrieval hit becomes an additional correct answer.
A value of 0 means that the retrieval gain leaves EM unchanged.
Benchmark ranks report only the denominator.
$T(R)$ reports how much of it reaches the answer for a fixed reader.

We pre-registered a reporting rule.
We report $T(R)$ only if $|\Delta\mathrm{Hit@5}| > 2$ percentage points (pp) and the Hit@5 difference is significant after Holm correction (\S\ref{sec:stats}).
Otherwise we mark $T(R)$ as not estimable, because a small or uncertain denominator makes the ratio unstable.
We compute its 95\% confidence interval with the paired bootstrap of \S\ref{sec:stats}: each resample of queries yields one ratio of mean differences, and we take the percentiles.

We also decompose EM by retrieval outcome:
\begin{equation}
\mathrm{EM} = h \cdot P(\mathrm{EM} \mid \mathrm{hit}) + (1 - h) \cdot P(\mathrm{EM} \mid \mathrm{miss}),
\label{eq:decomp}
\end{equation}
where $h$ is Hit@5.
This identity shows two routes by which a retrieval gain can fail to transfer.
The reader may answer incorrectly despite a hit, or it may answer correctly despite a miss.

\subsection{Corpus construction}
\label{sec:corpus}
Encoding the full BEIR-NQ corpus with six encoders exceeds our free GPU budget.
We therefore search a fixed 300,000-passage sub-corpus.
It contains all 4,201 judged relevant passages of the 3,452 judged test queries, 127,791 passages from the BM25 top 42 of those queries, and 168,008 random passages.
Every retrieval method searches the same sub-corpus.
This construction has a declared bias.
It contains BM25's own highest-ranked non-relevant passages, so it is harder for BM25 than for the dense methods.
It also removes most of the corpus, which raises absolute scores for all methods.
We therefore interpret relative comparisons among methods, not absolute scores.

\subsection{Retrieval methods}
\label{sec:retrievers}
We compare seven retrieval methods (Table~\ref{tab:T1}).
R0 is BM25 (bm25s 0.3.11, $k_1 = 1.2$, $b = 0.75$).
R1 to R6 are open-source dense encoders from 22.7M to 595.8M parameters: all-MiniLM-L6-v2, bge-small-en-v1.5, bge-base-en-v1.5, bge-large-en-v1.5, e5-base-v2 and Qwen3-Embedding-0.6B.
We pin each model to a fixed revision.
Each encoder truncates inputs to 256 tokens and runs in fp16.
We L2-normalise embeddings and search with exact inner-product search (FAISS flat index), which equals cosine similarity.
Exact search removes approximation error from the comparison.

\subsection{Reader}
\label{sec:reader}
The reader is Qwen2.5-7B-Instruct in 4-bit AWQ format, served with vLLM 0.30.0 on one Tesla T4.
It decodes greedily, with at most 32 new tokens, and stops at the first newline.
It receives the top $k = 5$ passages, each with its title, in rank order (prompt A1, Appendix~\ref{app:prompts}).
We run nine conditions on every query.
Seven use R0 to R6.
Closed-book (CB) gives no passages (prompt A2).
Gold gives only the judged relevant passages.
CB and Gold serve as reference points without and with the judged relevant passages.

\subsection{Judge}
\label{sec:judge}
The judge is Mistral-7B-Instruct-v0.3.
It receives exactly the passages the reader saw, the question and the answer, and it outputs SUPPORTED or NOT\_SUPPORTED (zero-shot prompt, Appendix~\ref{app:judge}).
We parse the verdict with a fixed regular expression and do not retry failures.
It judges the first 500 queries of our seeded query order for each retrieval method and for Gold.
The judge is unvalidated.
We compare it with human labels on a stratified sample of 200 items and report its scores under a pre-registered agreement gate (\S\ref{sec:results}).

\subsection{Statistics}
\label{sec:stats}
All comparisons are paired, because every method answers the same queries.
We test binary measures (Hit@$k$, EM, SUPPORTED) with the exact McNemar test.
We test graded measures (nDCG@10, token F1) with the Wilcoxon signed-rank test.
We correct each metric's 21 pairwise tests across the seven retrieval methods with the Holm procedure.
We report 95\% confidence intervals from a percentile bootstrap over queries, with 10,000 resamples and seed 0.
Before the main run, we estimated the pairwise EM discordance rate on a 200-query pilot as 0.165.
At this rate, with 2,255 queries and 80\% power, the minimum detectable EM difference is 3.4 pp at the worst-case Holm level ($\alpha = 0.05/21$).
We therefore report smaller EM differences as ``no difference detected'', not as equality.

\subsection{Stages}
\label{sec:stages}
The experiment runs in four stages (Figure~\ref{fig:F2}).
S1 builds the sub-corpus, encodes it with each method, retrieves the top 100 passages for all 3,452 judged queries, and measures cost: encoding throughput, corpus encoding time, query latency and peak GPU memory.
S3a runs the reader on the nine conditions for the 2,255 joined queries.
S3b runs three ablations: EM at $k \in \{1, 3, 10\}$ for R0, R1 and R4; an abstention prompt (A3) on 200 queries; and a determinism check that repeats 100 generations in a new session.
S4 runs the judge.
We fixed the design, metrics and tests in a pre-registration before S1 and record every deviation (\S\ref{sec:setup}).

\paragraph{Threats to validity inside the method.}
Each encoder uses the query and passage prefixes from its model card, for example ``query:'' and ``passage:'' for e5-base-v2.
The comparison is therefore between models as their authors intend them to be used, not a prefix-free comparison.
The sub-corpus favours dense methods over BM25 (\S\ref{sec:corpus}).
The reader, dataset and hardware are single choices, and our conclusions hold within them.
